\section{Taust ja kirjanduse ülevaade} \label{taust}

Käesolev peatükk annab ülevaate suurte keelemudelite arengust, eesti keele töötlemise spetsiifikast ning varasematest uurimustest tehisintellekti kasutamisest akadeemilises tekstiloomes.

\subsection{Suurte keelemudelite areng} \label{subsec:llm-areng}

Suurte keelemudelite ajalugu ulatub tagasi 2017. aastasse, mil Vaswani et al. avaldasid transformeri arhitektuuri \cite{vaswani2017attention}. See murdepunkt pani aluse kõigile kaasaegsetele keelemudelitele. Esimesed märkimisväärsed keelemudelid, nagu BERT \cite{devlin2019bert} ja GPT-2 \cite{radford2019language}, demonstreerisid, et suuremahuline eeltreenimine võimaldab mudelitel omandada laialdasi keelelisi teadmisi.

GPT-3 ilmumine 2020. aastal \cite{brown2020language} tõi kaasa paradigmamuutuse, näidates, et 175 miljardi parameetriga mudel suudab täita mitmesuguseid ülesandeid ilma spetsiifilise peenhäälestamiseta (\textit{fine-tuning}). Sellest ajast on mudelid kiiresti arenenud nii suuruse, arhitektuuri kui ka võimekuse poolest.

\subsection{Eesti keele töötlemine keelemudelites} \label{subsec:eesti-keel}

Eesti keel kuulub soome-ugri keelkonda ja on morfoloogiliselt rikas keel 14 käändega ning keerulise sõnamoodustuse süsteemiga \cite{erelt2007eesti}. See eripära teeb eesti keele töötlemise keelemudelite jaoks oluliselt keerulisemaks kui näiteks inglise keele puhul.

Suurte keelemudelite eesti keele oskus sõltub muu hulgas treeningandmete mahust ja kvaliteedist. Common Crawli 2026. aasta statistika järgi moodustas eesti keel HTML-lehtede põhikeelena viimastes crawl'ides ligikaudu 0,14\%, samal ajal kui inglise keele osakaal oli üle 40\% \cite{commoncrawl2026languages}. LLM-ide treeninguks puhastatud CulturaX andmestikus on vahe samas suurusjärgus: eesti keelt on 8,8 miljardit tokenit ehk 0,14\%, inglise keelt aga 2,85 triljonit tokenit ehk 45,13\% \cite{nguyen2024culturax}.

Tartu Ülikooli keeletehnoloogia uurimisrühm on teinud märkimisväärset tööd eesti keele ressursside ja tööriistade arendamisel. EstBERT \cite{tanvir2021estbert} on eesti keelele kohandatud BERT mudel, mis on treenitud eestikeelsel tekstikorpusel. Hiljutine EstLLM projekt \cite{dorkin2026estllm} on teinud olulisi edusamme suurte keelemudelite eesti keele võimekuse parandamisel, kasutades jätkutreenimist (\textit{continued pretraining}) ja järeltreenimist (\textit{post-training}) mitmekeelsete mudelite baasil. Uurimus näitas, et spetsialiseeritud treeninguga saab eesti keele kvaliteeti oluliselt tõsta, ilma et mudeli üldine arutlusvõimekus kannataks.

Mitmekeelsete mudelite arendamise seisukohast väärib esiletõstmist Aya mudel \cite{ustun2024aya}, mis on juhendipõhiselt peenhäälestatud avatud lähtekoodiga keelemudel, mis toetab 101 keelt, sealhulgas mitmeid väikese ressursiga keeli. Aya ületas mT0 ja BLOOMZ mudeleid enamikus ülesannetes, demonstreerides, et sihipärase mitmekeelse peenhäälestamisega on võimalik oluliselt parandada väikeste keelte kvaliteeti. Kuulmets et al. \cite{kuulmets2025finno} süstemaatiline hindamine näitas, et soome-ugri keelte, sealhulgas eesti keele, tugi avatud keelemudelites on viimastel aastatel kiiresti paranenud, kuid jääb endiselt inglise keelest oluliselt maha. Suurte keelemudelite puhul on ka eesti keele kvaliteet viimase kahe aasta jooksul märgatavalt paranenud. Kui varasemad GPT-3-põhised mudelid tegid eesti keeles sageli lihtsalt märgatavaid grammatilisi vigu, siis GPT-5.4, Gemini 3.1 Pro, Claude Sonnet 4.6 ja Claude Opus 4.6 suudavad luua grammatiliselt suures osas korrektset eesti keelt. Siiski esineb endiselt probleeme tehnilise terminoloogiaga, akadeemilise stiiliga ning pikematel tekstidel keelelise järjepidevuse hoidmisega. Kiil \cite{kiil2025bioloogia} on Tartu Ülikooli bakalaureusetöös võrrelnud suuri keelemudeleid eesti bioloogiaolümpiaadi küsimuste põhjal. Parim mudel saavutas 85,35\% täpsuse, mis on võrreldav olümpiaadiosalejate tipptulemusega. See kinnitab mudelite kasvavat võimekust ka eestikeelsetes ainespetsiifilistes ülesannetes.

\subsection{TI akadeemilises kirjutamises} \label{subsec:ai-akadeemia}

Tehisintellekti kasutamine akadeemilises kirjutamises on muutunud väga aktuaalseks teemaks. Liang et al. \cite{liang2024mapping} leidsid, et TI-genereeritud teksti osakaal teadusartiklites on märkimisväärselt kasvanud, eriti arvutiteaduse ja meditsiini valdkonnas. See tõstatab küsimusi akadeemilise aususe, originaalsuse ja kvaliteedikontrolli kohta.

Mitmed ülikoolid, sealhulgas Tartu Ülikool, on kehtestanud juhised TI kasutamiseks akadeemilises töös \cite{tuati2026juhend}. Üldiselt on lubatud kasutada TI-d abivahendina, kuid üliõpilane peab selgelt märkima, millises ulatuses ja millisel eesmärgil TI kasutati. Lõputöö ei tohi olla sisuliselt TI-genereeritud, vaid üliõpilane peab demonstreerima iseseisvat analüütilist mõtlemist.

Rodafinos \cite{rodafinos2025integration} rõhutab, et generatiivsete TI tööriistade vastutustundlik integreerimine akadeemilisse kirjutamisse on hädavajalik, tuues esile nii efektiivsuse kasvu kui ka kvaliteedi paranemise. Rõhutades samas TI väljundite kontrollimise ja inimese järelevalve säilitamise vajadusest. Lund et al. \cite{lund2025perceptions} leidsid 401 USA üliõpilase küsitluses, et tudengite eetilised veendumused, mitte institutsionaalsed poliitikad, on tugevaim ennustaja TI kasutamise tajumisel akadeemilise väärkäitumisena.

Varasemad uurimused on näidanud, et TI mudelid on eriti kasulikud järgmistes akadeemilise kirjutamise aspektides:
\begin{itemize}
    \item \textbf{Teksti struktureerimine} -- mudelid aitavad koostada tekstikavandeid, peatükkide ülesehitust ja kavandipõhist mustandit; Lee et al. \cite{lee2025navigating} näitasid, et eksplitsiitsed kavandid parandavad eri LLM-ide loodud tekstide kvaliteeti.
    \item \textbf{Kirjanduse ülevaade} -- mudelid suudavad toetada kirjanduse kogumist, organiseerimist, kokkuvõtmist ja sünteesimist, kuigi allikate täpsus vajab kontrollimist \cite{tang2025large,usher2025prompt}.
    \item \textbf{Keeleline toimetamine} -- ChatGPT, Claude'i, Gemini ja Copiloti tüüpi tööriistad võivad aidata grammatika, sõnastuse, sidususe ja loetavuse parandamisel, kuid nende väljund ei ole inimtoimetajaga samaväärselt usaldusväärne \cite{nadra2026genai}.
    \item \textbf{Tõlkimine} -- GPT-, Claude'i, Gemini ja DeepSeeki tüüpi mudelid suudavad pakkuda kasutatava kvaliteediga tõlkeid akadeemilisest ja tehnilisest tekstist, kuid valdkonnaspetsiifiline järeltoimetamine jääb vajalikuks \cite{pinto2026translation}.
\end{itemize}

Samas on tuvastatud olulisi puudujääke:
\begin{itemize}
    \item \textbf{Hallutsinatsioonid} -- mudelid võivad genereerida faktiliselt valesid väiteid või olematuid viiteid. Mugaanyi et al. \cite{mugaanyi2024citation} leidsid, et ChatGPT 3.5 genereeritud viidetest olid 72,7\% loodusteadustes ja 76,6\% humanitaarteadustes küll olemasolevad, kuid DOI-de ehk digitaalsete objektiidentifikaatorite (\textit{Digital Object Identifier}) täpsus oli vaid 32,7\% ja 8,5\% vastavalt. Uuemad hindamised näitavad, et probleem püsib ka kirjanduse ülevaate koostamise ja veebipõhiste agentsete töövoogude puhul: Tang et al. \cite{tang2025large} tuvastasid hallutsineeritud viiteid ka arenenud mudelite loodud kirjandusülevaadetes ning Rao ja Callison-Burch \cite{rao2026bibtex} hindasid GPT-5, Claude Sonnet 4.6 ja Gemini 3 Flashi veebipõhiseid väljundeid ning leidsid, et loodud BibTeX-kirjetest olid täielikult korrektsed vaid 50,9\% (üldine täpsus oli 83,6\%). Kim et al. \cite{kim2025geographic} tuvastasid lisaks, et hallutsinatsioonimäärad on geograafiliselt ebavõrdsed -- madalama sissetulekuga riikide kontekstis ületas DOI hallutsinatsioonimäär 80\%.
    \item \textbf{Sisuline sügavus} -- TI-genereeritud tekstid kipuvad olema pealiskaudsed ja väga üldistavad. Peters ja Chin-Yee \cite{peters2025generalization} näitasid 4900 teaduskokkuvõtte põhjal, et ChatGPT-4o, ChatGPT-4.5, DeepSeek, LLaMA 3.3 70B ja Claude 3.7 Sonnet kaldusid teadustulemusi üleüldistama; DeepSeek, ChatGPT-4o ja LLaMA 3.3 70B puhul esines seda 26--73\% juhtudest.
    \item \textbf{Originaalsus} -- keelemudelid ei loo iseseisvalt uusi teadmisi, vaid kombineerivad olemasolevat informatsiooni. Doshi ja Hauser \cite{doshi2024generative} leidsid, et generatiivne TI võib suurendada üksiku teksti loomingulisust, kuid vähendama samal ajal tekstikorpuse mitmekesisust. See toetab varasemat kriitikat mudelite tuletusliku olemuse kohta \cite{bender2021dangers}.
    \item \textbf{Väikeste keelte spetsiifika} -- eesti keele ja teiste väikeste keelte puhul on kvaliteet ebaühtlasem kui inglise keeles. Kuulmets et al. \cite{kuulmets2025finno} näitasid, et soome-ugri keelte tugi avatud keelemudelites on paranenud, kuid jääb suurtes keeltest korralikult maha. Lillepalu ja Alumäe \cite{lillepalu2025estonianbenchmark} rõhutavad emakeelsete eesti keele testide vajalikkust ning EstLLM projekt \cite{dorkin2026estllm} näitab, et eesti keele kvaliteet paraneb märgatavalt alles sihipärase jätkutreenimise ja järeltreenimisega.
\end{itemize}
